\section{What Sessions Keeps}
\label{sec:ch09-what-sessions-keeps}

\index{stub!what is left when a type moves out|(}%
The Sessions service still has to name \code{Speaker}, because
\code{Session.speaker} still returns one. What it no longer has is anything to
put in it. \code{Speaker.cs} loses two properties and gains an attribute:

\begin{minted}[highlightlines={14}]{csharp}
using HotChocolate.Types.Relay;

namespace Sessions;

/// <summary>
/// A speaker, as far as this service is concerned.
/// </summary>
/// <param name="Id">Stable identifier for this speaker.</param>
// An identifier and nothing else. The name and the bio moved to the
// Speakers service in chapter 9, so there is no row behind this record
// and no table to load it from. It is worth declaring anyway, because
// a client asking a session for its speaker's name needs somewhere to
// ask, and this type is that somewhere.
public sealed record Speaker([property: ID<Speaker>] int Id);
\end{minted}

\index{ID@\code{[ID<T>]}!makes the key a node id on the stub}%
\code{[ID<Speaker>]} is doing the work the \code{[NodeResolver]} used to do.
This type is not a \code{Node} here any more, so nothing else would make its
\code{id} export as \code{ID!} and carry the encoded form, and both of those
have to be true: the string this service hands out is the key the Speakers
service will be asked to resolve, so it must be the string the Speakers
service would have produced.

\index{stub!the type class shrinks to one line}%
\code{SpeakerType} loses both resolvers and most of itself:

\begin{minted}{csharp}
using HotChocolate.ApolloFederation.Types;

namespace Sessions;

/// <summary>
/// The stub. This service names Speaker so that Session.speaker has a
/// type, declares which field identifies one, and then says it cannot
/// turn such a key back into anything: resolvable false is the whole
/// difference between this and the Speakers service's declaration of
/// the same type. Without it, composition would believe two subgraphs
/// can both answer for a speaker and route accordingly.
/// </summary>
[ObjectType<Speaker>]
[Key("id", false)]
public static partial class SpeakerType;
\end{minted}

\index{Key@\code{[Key]}!its second argument}%
That second argument is the chapter's first new piece of federation.
\code{[Key]} has exactly two constructors: one taking nothing, and one taking
a field set and a boolean that defaults to true. Set it false and the exported
schema says so.

\index{resolvable@\code{resolvable: false}!what Apollo calls a stub}%
Chapter~\ref{ch:federation-model} read that argument out of the directive
definition and said what setting it false declares. Apollo has a name for the
type you end up with: a \emph{stub
definition}~\autocite{apollocontributefields}, carrying the key fields of an
entity and nothing else. That is this file exactly, and it is worth having the
word for it, because a stub is a thing you will write once per seam for the
rest of your graph's life.

\subsection{The Resolver That Does No Work}

\index{Session.speaker@\code{Session.speaker}!costs no statement now}%
\code{Session.speaker} keeps its resolver and loses its body:

\begin{minted}[highlightlines={9,10,11}]{csharp}
    /// <summary>
    /// The person giving this session.
    /// </summary>
    // Since chapter 9 this service holds the identifier and nothing
    // else, so the whole of the work is wrapping the foreign key it
    // already has in the type the schema promises. No statement is
    // issued and nothing is checked: if no speaker carries this
    // number, the Speakers service is where that is discovered.
    public static Speaker? GetSpeaker(
        [Parent(nameof(Session.SpeakerId))] Session session)
        => new(session.SpeakerId);
\end{minted}

\index{DataLoader!the one this service no longer needs}%
No \code{async}, no DataLoader, no cancellation token. The method is now a
constructor call over a column the row already carried, and
\code{SpeakerDataLoader.cs} is deleted from this service, having moved to the
one with the table in it.

\index{Parent@\code{[Parent]}!still load-bearing, and now silently}%
The \code{[Parent]} attribute naming \code{SpeakerId} survives, and it got
more dangerous rather than less.
Chapter~\ref{ch:data-without-n-plus-1} put it there so the
projection would keep \code{SpeakerId} in the row when nobody selected it, and
measured what happens without it: every speaker resolves to null. Drop it now
and the failure is quieter. The column comes back as the default for its type,
the record is constructed from a zero, and every session reports a speaker
whose node id decodes to \code{Speaker:0}. That is a well-formed key for a row
that does not exist, and the Speakers service will answer it with a null and
no error at all.

\index{nullability!ownership stops being a judgment call}%
\code{speaker} stays nullable, and the reason is better than it was.
Chapter~\ref{ch:schema-design} left it nullable as a judgment about what this
service could guarantee, and made it stick by measuring what
\code{speaker: Speaker!} does to a response when one row is missing. It is not
a judgment any more. This service cannot guarantee a speaker because the
speakers are in another process, and section~\ref{sec:ch09-the-first-query-across-two}
shows what arrives when that process is not running.

\subsection{Three Readouts of One Argument}

\index{schema!the stub as exported}%
Export the schema and the type is four lines and its description:

\begin{graphqlsdl}
"A speaker, as far as this service is concerned."
type Speaker @key(fields: "id", resolvable: false) {
  "Stable identifier for this speaker."
  id: ID!
}
\end{graphqlsdl}

\index{Node@\code{Node}!the stub does not implement it}%
\code{implements Node} is gone, because there is no node resolver to make it
true. So is every field except the key. A reader of this document learns that
a speaker exists, that it is identified by \code{id}, and nothing else, which
is exactly as much as this service knows.

\index{Entity@\code{\_Entity}!shrinks by one}%
The second readout is the union chapter~\ref{ch:first-subgraph} called a
readout:

\begin{graphqlsdl}
union _Entity = Session
\end{graphqlsdl}

\index{Entity@\code{\_Entity}!is what resolvable false changes}%
Two members at the end of chapter~\ref{ch:first-subgraph}, one now. That union
is generated from the types this service will answer an entity fetch for, and
\code{resolvable: false} is what took \code{Speaker} out of it. If you make a
stub and the type is still in the union, the argument did not take.

\index{router.json@\code{router.json}!compiles resolvable false to a flag}%
The third readout is in the router execution config, and it is the one that
does the routing. Chapter~\ref{ch:composition} read the \code{@key} directives
out of that file as a lookup table and called it the seam in the form the
router consumes. Compose the graph now and the Sessions entry reads:

\begin{minted}{json}
"keys": [
  {"typeName": "Session", "selectionSet": "id"},
  {"typeName": "Speaker", "selectionSet": "id",
   "disableEntityResolver": true}
]
\end{minted}

\index{disableEntityResolver@\code{disableEntityResolver}!is the whole mechanism}%
The stub is in the key table. It has to be: the router still needs to know the
shape of a speaker's key, because \code{Session.speaker} hands one back and
something has to be sent onward. What it also has is one boolean, and that
boolean is the entire mechanism. The Speakers subgraph's entry for the same
type carries no such property, because the flag is written only when it is
false, so its absence there is the positive case.

\index{resolvable@\code{resolvable: false}!is advice to the composer, not a guard}%
One thing that flag is not is a guard. It tells the composer where not to
route, and nothing enforces it on the service itself. Send the Sessions
subgraph a representation naming a \code{Speaker}, on its own port, and it
does not answer with a polite null:

\begin{minted}[fontsize=\footnotesize,breakanywhere]{json}
{"errors":[{"message":"Unexpected Execution Error","path":["_entities"]}],"data":null}
\end{minted}

\index{entities@\code{\_entities}!the stub is a hole in it}%
Not a null in the list, which is what chapter~\ref{ch:first-subgraph} measured
for a key matching no row. An unhandled error, taking the whole \code{data}
key with it. Nothing a client can see says this would happen: the type is in
the schema, it advertises a key, and \code{resolvable: false} lives in a
directive that ordinary introspection does not show. Chapter~\ref{ch:entities-authorization}
is about what the entity route hands over when asked politely, and this is one
more thing it will find waiting.
\index{stub!what is left when a type moves out|)}%
